\documentclass[11pt]{article}
\usepackage{amsmath,amssymb}

\setlength{\parindent}{0pt}
\setlength{\parskip}{0.7em}
\setlength{\textwidth}{6.5in}
\setlength{\oddsidemargin}{0in}
\setlength{\evensidemargin}{0in}
\setlength{\topmargin}{-0.4in}
\setlength{\textheight}{9.2in}

\begin{document}

\begin{center}
{\Large SYSC 5108 Exam Study: Assignment 3, Question 5}\\
\vspace{0.3em}
{\large Forward pass and backpropagation in a 1D CNN with max pooling}\\
\vspace{0.3em}
Source question: \texttt{SYSC 5108 W26 Assignment 3.pdf}, Question 5
\end{center}

\section*{Question Summary}

The network is a one-dimensional convolutional neural network with:
\begin{itemize}
\item a shared ReLU filter applied at neurons 1, 2, 3, and 4,
\item two max-pooling neurons 5 and 6, each pooling two neighboring activations with no overlap,
\item one fully connected ReLU output neuron 7.
\end{itemize}

The input is
\[
x=[1,1,1,0,0,0].
\]
The shared filter has bias and weights
\[
w_0=0.75,\qquad w_1=1,\qquad w_2=1,\qquad w_3=-1.
\]
The filter width is 3 and the stride is 1.

The final output neuron has
\[
w_{7,0}=0.1,\qquad w_{7,5}=0.2,\qquad w_{7,6}=0.5.
\]
All neurons use ReLU:
\[
\operatorname{ReLU}(z)=\max(0,z).
\]
The target output is
\[
t=1.
\]

\section*{Step 1: Relate the Problem to the Course Material}

Chapter 5 slides define convolutional neural networks using three core ideas:
\begin{itemize}
\item \textbf{Sparse interactions}: each neuron sees only a local window of the input.
\item \textbf{Parameter sharing}: the same filter weights are reused across different locations.
\item \textbf{Pooling}: nearby outputs are summarized, often by max pooling.
\end{itemize}

The slides also note that each kernel creates one feature map and that pooling produces a summary statistic of nearby outputs. That is exactly what happens in this problem.

\section*{Part (a): Compute the Network Output}

\subsection*{1. Convolution windows}

With filter width 3 and stride 1, the four windows are:
\[
[1,1,1],\qquad [1,1,0],\qquad [1,0,0],\qquad [0,0,0].
\]

Each convolution neuron computes
\[
z=w_0+w_1x_1+w_2x_2+w_3x_3,
\qquad
a=\operatorname{ReLU}(z).
\]

\subsection*{2. Convolution layer activations}

Neuron 1 sees $[1,1,1]$:
\[
z_1=0.75+(1)(1)+(1)(1)+(-1)(1)=1.75.
\]
So
\[
a_1=\operatorname{ReLU}(1.75)=1.75.
\]

Neuron 2 sees $[1,1,0]$:
\[
z_2=0.75+(1)(1)+(1)(1)+(-1)(0)=2.75.
\]
So
\[
a_2=2.75.
\]

Neuron 3 sees $[1,0,0]$:
\[
z_3=0.75+(1)(1)+(1)(0)+(-1)(0)=1.75.
\]
So
\[
a_3=1.75.
\]

Neuron 4 sees $[0,0,0]$:
\[
z_4=0.75+(1)(0)+(1)(0)+(-1)(0)=0.75.
\]
So
\[
a_4=0.75.
\]

Therefore the feature map is
\[
[a_1,a_2,a_3,a_4]=[1.75,2.75,1.75,0.75].
\]

\subsection*{3. Max-pooling layer}

Neuron 5 pools neurons 1 and 2:
\[
a_5=\max(1.75,2.75)=2.75.
\]

Neuron 6 pools neurons 3 and 4:
\[
a_6=\max(1.75,0.75)=1.75.
\]

\subsection*{4. Output neuron}

Neuron 7 computes
\[
z_7=w_{7,0}+w_{7,5}a_5+w_{7,6}a_6.
\]
Substitute:
\[
z_7=0.1+(0.2)(2.75)+(0.5)(1.75).
\]
So
\[
z_7=0.1+0.55+0.875=1.525.
\]
Apply ReLU:
\[
a_7=\operatorname{ReLU}(1.525)=1.525.
\]
Therefore the network output is
\[
\boxed{1.525.}
\]

\section*{Part (b): Compute the Delta Values}

Use squared error:
\[
E=\frac12(y-t)^2,
\qquad y=a_7.
\]
Since
\[
y=1.525,\qquad t=1,
\]
the output error signal is
\[
\delta_7=\frac{\partial E}{\partial z_7}
=(y-t)\operatorname{ReLU}'(z_7).
\]
Because $z_7=1.525>0$,
\[
\operatorname{ReLU}'(z_7)=1.
\]
So
\[
\delta_7=(1.525-1)(1)=0.525.
\]

\subsection*{1. Backpropagate to pooling neurons}

Neuron 5 connects to neuron 7 with weight $w_{7,5}=0.2$, so
\[
\delta_5=w_{7,5}\delta_7=(0.2)(0.525)=0.105.
\]

Neuron 6 connects to neuron 7 with weight $w_{7,6}=0.5$, so
\[
\delta_6=w_{7,6}\delta_7=(0.5)(0.525)=0.2625.
\]

\subsection*{2. Route max-pooling gradients}

Max pooling sends the gradient only to the selected maximum.

Neuron 5 chose neuron 2 because
\[
\max(a_1,a_2)=\max(1.75,2.75)=a_2.
\]
Therefore
\[
\delta_2=0.105,\qquad \delta_1=0.
\]

Neuron 6 chose neuron 3 because
\[
\max(a_3,a_4)=\max(1.75,0.75)=a_3.
\]
Therefore
\[
\delta_3=0.2625,\qquad \delta_4=0.
\]

\subsection*{3. ReLU derivatives at convolution neurons}

All convolution pre-activations are positive:
\[
z_1=1.75,\quad z_2=2.75,\quad z_3=1.75,\quad z_4=0.75.
\]
So every ReLU derivative is 1, and the deltas stay unchanged.

Hence the full set of deltas is
\[
\boxed{
\delta_1=0,\quad
\delta_2=0.105,\quad
\delta_3=0.2625,\quad
\delta_4=0,\quad
\delta_5=0.105,\quad
\delta_6=0.2625,\quad
\delta_7=0.525.
}
\]

\section*{Part (c): Calculate the Shared Filter Weight Updates}

Because the filter is shared, the gradient for each filter parameter is the sum of its contribution from all four convolution neurons.

\subsection*{1. Bias term \(w_0\)}

Each convolution neuron uses the same bias weight, so
\[
\frac{\partial E}{\partial w_0}
=
\delta_1+\delta_2+\delta_3+\delta_4.
\]
Substitute:
\[
\frac{\partial E}{\partial w_0}
=0+0.105+0.2625+0=0.3675.
\]

\subsection*{2. Weight \(w_1\)}

\(w_1\) multiplies the first entry of each window. The first entries are
\[
1,\ 1,\ 1,\ 0.
\]
So
\[
\frac{\partial E}{\partial w_1}
=
\delta_1(1)+\delta_2(1)+\delta_3(1)+\delta_4(0).
\]
Thus
\[
\frac{\partial E}{\partial w_1}
=0+0.105+0.2625+0=0.3675.
\]

\subsection*{3. Weight \(w_2\)}

\(w_2\) multiplies the second entry of each window. Those entries are
\[
1,\ 1,\ 0,\ 0.
\]
So
\[
\frac{\partial E}{\partial w_2}
=
\delta_1(1)+\delta_2(1)+\delta_3(0)+\delta_4(0).
\]
Thus
\[
\frac{\partial E}{\partial w_2}=0+0.105+0+0=0.105.
\]

\subsection*{4. Weight \(w_3\)}

\(w_3\) multiplies the third entry of each window. Those entries are
\[
1,\ 0,\ 0,\ 0.
\]
So
\[
\frac{\partial E}{\partial w_3}
=
\delta_1(1)+\delta_2(0)+\delta_3(0)+\delta_4(0).
\]
Thus
\[
\frac{\partial E}{\partial w_3}=0.
\]

Therefore the shared-filter gradients are
\[
\boxed{
\frac{\partial E}{\partial w_0}=0.3675,\qquad
\frac{\partial E}{\partial w_1}=0.3675,\qquad
\frac{\partial E}{\partial w_2}=0.105,\qquad
\frac{\partial E}{\partial w_3}=0.
}
\]

If the learning rate is \(\alpha\), then gradient descent gives
\[
\Delta w_0=-\alpha(0.3675),\qquad
\Delta w_1=-\alpha(0.3675),\qquad
\Delta w_2=-\alpha(0.105),\qquad
\Delta w_3=0.
\]
Equivalently,
\[
w_0^{\text{new}}=w_0-\alpha(0.3675),\quad
w_1^{\text{new}}=w_1-\alpha(0.3675),\quad
w_2^{\text{new}}=w_2-\alpha(0.105),\quad
w_3^{\text{new}}=w_3.
\]

\section*{Why This Problem Matters}

This question captures three CNN ideas that are easy to miss on an exam:
\begin{itemize}
\item \textbf{Convolution reuses one filter} at multiple positions, so gradients for shared weights must be added across all positions.
\item \textbf{Max pooling is not linear}. During backpropagation, only the input that achieved the maximum receives the upstream gradient.
\item \textbf{ReLU derivatives are simple}. If the pre-activation is positive, the derivative is 1; otherwise it is 0.
\end{itemize}

\section*{Concise Exam Answer}

Convolution windows:
\[
[1,1,1],\quad [1,1,0],\quad [1,0,0],\quad [0,0,0].
\]
Convolution outputs:
\[
z_1=1.75,\quad z_2=2.75,\quad z_3=1.75,\quad z_4=0.75,
\]
so
\[
a_1=1.75,\quad a_2=2.75,\quad a_3=1.75,\quad a_4=0.75.
\]
Pooling gives
\[
a_5=\max(1.75,2.75)=2.75,\qquad
a_6=\max(1.75,0.75)=1.75.
\]
Output neuron:
\[
z_7=0.1+0.2(2.75)+0.5(1.75)=1.525,
\qquad
a_7=1.525.
\]
Thus the output is
\[
\boxed{1.525.}
\]
With target \(t=1\),
\[
\delta_7=(1.525-1)=0.525,
\quad
\delta_5=0.2(0.525)=0.105,
\quad
\delta_6=0.5(0.525)=0.2625.
\]
Pooling routes these to the maxima:
\[
\delta_2=0.105,\quad \delta_3=0.2625,\quad \delta_1=\delta_4=0.
\]
Therefore
\[
\boxed{
\delta_1=0,\ \delta_2=0.105,\ \delta_3=0.2625,\ \delta_4=0,\ \delta_5=0.105,\ \delta_6=0.2625,\ \delta_7=0.525.
}
\]
Shared-filter gradients:
\[
\boxed{
\frac{\partial E}{\partial w_0}=0.3675,\quad
\frac{\partial E}{\partial w_1}=0.3675,\quad
\frac{\partial E}{\partial w_2}=0.105,\quad
\frac{\partial E}{\partial w_3}=0.
}
\]

\section*{Cross References Used}

\begin{itemize}
\item Assignment: \texttt{SYSC 5108 W26 Assignment 3.pdf}, Question 5.
\item Local submitted Assignment 3 PDF checked for comparison, Question 5 response.
\item Slides: Chapter 5 slide deck on convolutional neural networks, convolution, stride, pooling, feature maps, and trainable parameters.
\item Slides: Chapter 4 slide deck on ReLU activation functions.
\item Textbook: Goodfellow, Bengio, and Courville, \textit{Deep Learning}, Chapter 6 on feedforward networks and ReLU activations.
\item Textbook: Goodfellow, Bengio, and Courville, \textit{Deep Learning}, Chapter 9 on convolution, parameter sharing, and pooling concepts.
\end{itemize}

\end{document}
